\chapter{The patterns we experience, the laws we accept}
\label{Intro_Physics}

\epigraph{The good thing about science is that it's true whether or not you believe in it}{\textit{ Neil DeGrasse Tyson}}

\section{Science doesn't know everything!}

\noindent We observe patterns of reality. The sun rises every morning. So we postulate that it will rise tomorrow too. And the day after. And the day after that. This pattern does not seem to change. Though the details of exactly at what time this happens, how long it takes for the sun to rise from over the horizon to its zenith, the angle under which it does so, the heat it radiates... do vary when looked at closely. But despite these differences there is a clear pattern to discern. We are convinced the sun will rise tomorrow. It is not impossible we are wrong about this. But all the evidence at hand does not contradict this prediction, it corroborates it.\\

\noindent Science does not know everything. Far from it. But to quote the comedian/physicist Dara O'Briain: "Science knows it doesn't know everything. Otherwise, it'd stop!".\\

\noindent Many claims have been made and tested. Some turned out to be in contradiction with what actually happens in reality. We know that these claims that have been falsified are untrue. There is therefore quite a bit that science does know, insofar as one can know anything.\\ 

\noindent In this chapter I have two main goals:
\begin{enumerate}
\item Discuss a number of these topics about which we do know quite a lot. 
\item Introduce a number of ideas and topics in physics without diving into the hard quantitative descriptions just yet. I will mostly omit math in this chapter and will focus on the concepts. In later chapter we will revisit many of these topics and include the mathematical description of them.
\end{enumerate}

\section{Newton's laws}

\noindent Mechanics, the body of knowledge concerning motion: how do objects move and why do they move that way. Traditionally one of the first topics treated in physics. And the basics of mechanics are given by Newton's three laws of motion.

\begin{enumerate}

\item \textbf{Newton's first law:} Any object remains in its state of motion, unless it is acted upon by a force.\\

\noindent Objects that are moving will keep on moving with the same speed in the same direction forever, unless a force compels them otherwise. Objects that are standing still will remain still unless a force compels them otherwise. This phenomenon is called \textbf{inertia}. This is the first of several instances in which we use a word used in daily life but meaning something (slightly) different in physics. Be aware of these differences as they may be the source of confusion or misunderstanding.\\

\noindent We experience inertia in several everyday circumstances:
\begin{itemize}
\item Sitting in a car that suddenly brakes, we are thrown forwards. Without the seat belts we could be thrown straight out through the front shield. The car's breaks do stop the car but our body does not experience this and inertia keeps it moving forward.

\item Taking a turn with the car, we experience a force pushing us away from the direction towards which the car is turning. Inertia actually keeps us moving straight ahead, but as the car turns we experience this as if there were a force on us in the other direction. 

\item When standing on a bus we fall backwards when it starts moving. The engine forces the bus to move but our body remains where it was standing. We experience this as a force backward as seen from the bus.

\end{itemize}
 
\noindent So why does everything around seem to come to a standstill? Well, because the surface on which the object is moving, as well as the air around it exert a frictional force on the moving object which makes it stop. The first of Newton's law states that in the absence of such forces, once in motion the object remains indefinitely in motion. An astronaut drifting away from her space-shuttle will keep on drifting indefinitely.

\item \textbf{Newton's second law:} In the presence of forces the object's velocity changes. The mass of the object is a measure for the difficulty of changing its velocity.\\

\noindent The equation known as Newton's second law, describing the motion of an object is given in its simplest form by
\begin{equation}
F = m a
\end{equation}
where $F$ is the force exerted on the object, $m$ is its mass and $a$ its acceleration. The acceleration, as we will see in the near future, is nothing but the change of the velocity. Let us consider two objects at rest, one with mass = $2 kg$ and the other with mass $1\; 000 kg$. We exert an equal force on both objects, equal to 10 Newton\footnote{Just like the unit of mass is called kg, the unit of force is called Newton. We will see in the chapter on mechanics what a force of 1 Newton represents.}. Looking at their acceleration (i. e. we look at how fast they move after the force was applied) we notice the following:
\begin{center}
\begin{tabular}{|c|c|c|}
  \hline
	& & \\
  \textbf{Object} &2 kg  & 1\;000 kg \\
	\hline
  & & \\
	\textbf{Speed} & $5 m/s $ & $0,01 m/s $\\
  \hline
\end{tabular}
\end{center}
\noindent The lighter object's acceleration is greater than that of the heavier one under the influence of the same force.\\

\noindent Kicking a football lying on the ground will make it fly away. The same kick will not do anything noticeable to a train. Equal forces on vastly different masses have vastly different consequences.\\

\noindent The gravitational force that the earth exerts on a hanging apple is as strong as the gravitational force that the apple exerts on the earth. But it is clearly the apple falling towards the earth and not the other way around. Why? Because the earth is enormously heavier compared to apple. So the change of velocity of the apple is much larger than that of the earth.\\

\noindent Newton's second law contains much more information than I explain here. In the chapters on vectors and on mechanics we will deepen our understanding of the content of this law and its predicted physical consequences.\\

\item \textbf{Newton's third law:} Forces come in equal but opposite pairs.

\begin{itemize}
\item While ice-skating you push your friend away. Not only will you friend be pushed forward (i.e. away from you), but you yourself will be pushed backwards (i.e. away from your friend). You exert a force on her. And she consequently exerts a force on you of equal size and opposite direction to the original force.
\item A rocket forces fuel out of its tail. The fuel consequently exerts a equal force on the rocket upward, which causes it to accelerate upwards.
\end{itemize}

\end{enumerate}

\noindent When describing the motion of a football, a rocket, an elephant and more generally in "`everyday circumstances"' these three laws give an accurate representation of what actually happens. But only the first of the laws does not have any known exceptions. The second and third law are known not to correctly describe the very small, the very fast, the very heavy, the very warm, the very cold,...\\

\noindent In such cases other laws hold, some of which we will discuss further on in the text. So keep in mind that together with a physical theory come its limits of applicability. And using the theories outside these limits will result in false predictions about what reality will/may do under these circumstances.

\section{Unification: l'union fait la force!}

\noindent There are many known forces in physics: gravitational, van der Waals,  electrical, magnetic, nuclear, drag, lift, frictional, centripetal, centrifugal,...\\

\noindent However, it turns out that there are only four fundamental ones:
\begin{enumerate}
\item Gravity
\item Electromagnetism
\item Strong nuclear
\item Weak nuclear
\end{enumerate}
We will discuss each of them in due time. But for now, it is sufficient to know that we do not know any other forces. All the other forces like friction or the centripetal force are one or several of the fundamental forces in disguise. We will come back to this point later.\\

\noindent This section is about forces. But about their similarities rather than their differences.
\begin{itemize}
\item A magnet pushing another magnet away seems to be something very different from the spark you feel (and sometimes even see) when touching a car door on a dry day.
\item The needle of a compass reacting to the earth's magnetic field seems to be completely different from the lightning strike during a thunderstorm.
\item An apple falling seems to be unrelated to the way the our solar system moves within the milky way.
\end{itemize}
Remember that impressions can be misleading. In each of these examples the phenomena are not only related, but turn out to be a consequence of the same underlying mechanism. Realizing that seemingly different phenomena can be understood in the same way is known in physics as unification.\\

\noindent Isaac Newton himself realized that the fall of an apple and the motion of the heavenly bodies are both instances of gravitational forces. But it doesn't stop there. The tides of the sea, and the presence of an atmosphere on some planets and not on others as well as the chemical constitution of that atmosphere can be explained by the same phenomenon that makes an apple fall from the tree.

\noindent Michael Faraday, James Clerk Maxwell and other $19^{th}$ century scientists realized that electric and magnetic phenomena are very closely related. When electric phenomena are present, magnetic phenomena are rarely far away and vice versa. Albert Einstein understood a couple of decades later that electric and magnetic phenomena are in some sense the same phenomenon, that they are two sides of a same coin (see section \ref{ElectromagneticSection}). Sailing the seas in a thunderstorm you may notice that your compass loses all sense of direction and starts turning randomly around when a lightning bolt strikes not too far away. A credit card or a cell phone can become useless if it comes in contact with a strong magnet.\\

\noindent Unification is a theme that comes up in all of science, not only in physics. The mechanism that explains why children often look like their parents is the same as the one explaining that a banana, a dragonfly and we all share the same chemical building blocks and blueprint.\\

\noindent Finding out that seemingly different phenomena are related is among the deepest insights science has of reality.

\section{Some things are better left conserved}

\noindent We are born, we grow up, we get old,...\\
As times goes on everything changes. Well, not exactly everything. Some things do not change. When a quantity does not change we say that it is conserved. Scientists discovered/defined several quantities that do not change. As far as we can tell, the universe came to be with a certain amount of those quantities which remain what they have always been.

\subsection{Conservation of energy}

\noindent The most famous of these quantities is probably energy. It is a painfully misunderstood concept and I am going to take some time to explain what it is and what it most certainly is not.\\

\subsubsection{Kinetic energy}

\noindent Let us consider a billiard table as seen from above. A very special table. One without any holes. But what makes it even more special is that it is completely frictionless. It is standing in a special room completely devoid of any air (so that no air resistance can be present). On this table there is one single billiard ball. A black one. Standing still. This ball has a mass which we will call $m_1$. It's speed will be called $v_1$ and right now it is 0.\\

\noindent We now step towards the table with a cue stick and hit the ball one time in some direction (probably wearing an astronaut suit to be able to breath and survive in the room). The ball's speed is still called $v_1$ but it isn't 0 anymore.\\

\noindent We now leave the table with the ball to its own devices for a while. It can even be a long while. Eventually, let's say after 27 years, we come back to our table. And if in the course of all this time no one touched to ball, we will notice something remarkable: the black ball's speed is exactly the same as it was when we left it all these years ago. On a real billiard table this ball will roll for a while until it comes to a standstill. But on this very special table there isn't any friction around to cause resistance on the ball. Without anything to prohibit the ball's continuous motion it will just keep on rolling. Of course it will hit the cushions of the billiard table. But this will only change the direction in which it moves not its speed.\\

\noindent This thought experiment may or may not seem obvious to you. But the following most probably will not be. On this very same table we do a second experiment. There are now 2 balls on the table: a black one (ball nr. 1) and a red one (ball nr. 2). Their masses will be called $m_1$and $m_2$ and their speeds $v_1$ and $v_2$ both of which are initially 0. We once again come and hit the black ball with a cue stick. This causes it to move so that $v_1 \neq 0$ but $v_2$ still remains zero. We once again leave the table for a very long time.\\

\noindent Upon our return we find both balls in motion. This of course does not surprise us. We realize that the black ball probably collided with the red one at some point and it too started to move. Over all this time the balls most probably collided many times. Our observation that the black ball's speed is the same when we left the table and when we return is not true anymore. We wonder if we there is some relation between the the speed of the black ball when we left and the speed of both balls right now. Care to venture a guess?\\

\noindent To distinguish between the speeds before we left and after coming back we write $v_1$ and $v_2$ (which is 0) for the speeds right after hitting the black ball, and $v_1'$ and $v_2'$ for those when we return. A first guess you may venture is that the sum of both speeds remains constant: $v_1 + v_2 = v_1' + v_2'$. This is quite a simple equation and would allow us to calculate the speed of the second ball if we knew the initial speeds of both balls and the current speed of the other one. But this guess is wrong! It is not what actually happens. So another guess...\\

\noindent After many guesses that all fail to describe what we actually see on the table we may come up with the following very weird equation:
\begin{equation}
\label{Kin_en}
\frac{1}{2} m_1 \times v_1 \times v_1 + \frac{1}{2} m_2 \times v_2 \times v_2 = \frac{1}{2} m_1 \times v_1' \times v_1' + \frac{1}{2} m_2 \times v_2' \times v_2'   
\end{equation}
It is indeed a weird combination of the speeds and the masses of the balls. But miraculously it is correct. The world could have been different. But in our world this equality holds under these circumstances\footnote{It is true that such circumstances do not exist: there is no such thing as a completely frictionless table. We can nevertheless approximate it by taking real tables and making them as slippery as we can get them. What is observed is that the slipperier it becomes, the closer what actually happens resembles the prediction of this idealized situation.}.\\

\noindent In the past the quantity $\frac{1}{2} m v^2$ was known as the "`vis viva"' of the object (life force). Today that name is not in use anymore and it currently goes through life as "`kinetic energy"'. What equation (\ref{Kin_en}) says is that if we add the kinetic energy of the red ball to the that of the black ball, the total does not change: "`The total kinetic energy is conserved"'.

\subsubsection{Gravitational potential energy}
\label{Grav_pot_en_sect}
\noindent Hold a pen still in your hand (still in that same air-free room). Then suddenly let it go. The pen falls. Its kinetic energy while in your hand is 0 because its speed is 0. But once you let it go it start moving faster and faster. So its kinetic energy increases. Does that mean that the "`law"' that we just discovered about the conservation of the total kinetic energy is wrong?\\

\noindent Remember that we said that the story about the billiard table only holds if during our absence no one came and gave the ball another hit with a cue. That is, the story only holds if during all that time no external force was exerted on the balls. If someone did come in during that time and for instance just stopped both balls and left them at rest, then they will of course remain so. The point is, in my story above no external force was exerted on the balls. While in our example of the pen there is an external force: gravity. The kinetic energy indeed does change.\\

\noindent The pen's kinetic energy clearly increases during the fall. But it turns out that there is a quantity that does not change during the fall:
\begin{equation}
\frac{1}{2} m v_1^2 + 10 m h_1 = \frac{1}{2} m v_2^2 + 10 m h_2
\end{equation}
As the pen falls its speed increases and its height $h$ decreases in such a way that the quantity $\frac{1}{2} m v^2 + 10$ m h does not change. Once again I want to emphasize that there is no obvious a priori reason for the world to be this way. But it is. The term $10 m h$ that is added here is called the gravitational potential energy. For a falling pen the sum of kinetic and gravitational potential energy remains constant throughout the fall. The kinetic energy is not conserved nor is the potential energy. But the total energy (i.e. their sum) is.

\subsubsection{Other energy forms}

\noindent When the falling pen finally lands on the floor it comes to a halt. Its kinetic energy becomes 0 and so does its gravitational potential energy as its height is now 0. Does that mean that the total energy of the world changed? No!\\

\noindent When the pen was at its highest point, just after you let it go, it had some gravitational potential energy. During the fall this energy is changed into kinetic energy as the pen's speed increases. The impact on the ground transfers this energy carried by the pen into kinetic energy of the molecules in the floor. The floor will have "`absorbed"' the fall. The motion of the pen is transferred during the impact to the molecules of the floor. They start shaking and moving harder which, as we will see in section (\ref{Temperature_Intro}) is thermal energy. The floor becomes a little warmer. But the total energy of the world remains unchanged.\\

\noindent There are still other forms of energy than the ones mentioned above: nuclear, chemical, elastic potential, electromagnetic,...\\
As we learn more physics we will encounter some of these and try to understand what they represent. But the bottom line for all of them is their sum never changes.

\subsubsection{What is energy?}

\noindent The term energy is used in ordinary speech in many (often vague and metaphorical) ways, some of which make more sense than others. But in the physical sciences energy is a well-defined and very abstract quantity that has the amazing property that it does not change. The kinetic energy of an object is the square of its speed multiplied by its mass and divided by 2. The kinetic energy, just like any other form of energy does not have a shape, a color, a taste or any directly tangible property. It's the result of a calculation. We calculate all the energies that are currently present. And whether we do this calculation now, 700 years ago or in 15 centuries, the result  is always the same number.\\

\noindent The total energy is constituted of a plethora of energy forms as mentioned above. Each energy form separately need not be conserved. Amounts of energy can change from one form to another and from one carrier to another. Without the total ever changing.\\

\noindent One surprising and highly non-trivial fact about conservation of total energy in the universe is that it turns out to be related to the the constancy in time of the laws of physics themselves. As far as we know the laws of nature are today what they were in the past, both recent and distant. And we do not have any empirical arguments to suggest they will change in the future. It's possible to show mathematically that if the laws of nature themselves change with time, the total energy of the universe will also change. Though the mathematics and the physics of this argument are beyond what we can discuss here and now.

\subsection{Conservation of linear and angular momentum}

\noindent Energy is not the only quantity that is conserved. Two other important conserved quantities are the linear momentum and the angular momentum. They represent in some sense the amount of linear motion (movement in a straight line, without any turns or rotations) and angular momentum (turning motion). The actual mathematics of these will be discussed in the chapter on mechanics. But here is already a small taster.

\subsubsection{Linear momentum}

\noindent You are in space completely at rest with a basketball in your hand far away from any other stars, planets or any object. Whatever the exact definition is of linear momentum\footnote{I conform from now on to the convention of speaking simply of momentum when meaning the linear momentum and will keep on calling its rotational counterpart angular momentum.}, it seems reasonable that if you and the ball are at rest both your momenta are 0. Consequently, the total momentum is 0. Conservation of momentum means that the total momentum remains 0. You now throw the ball away from you. It starts moving away from you: it has momentum. The ball's momentum is not zero anymore. But if conservation of total momentum holds, you need to be moving in the opposite direction so that the total momentum of you and the ball remains 0. That is indeed what is observed in reality\footnote{On earth though this does not seem to hold. When you kick a football away, the ball flies off but you do not fly backwards to compensate compensate. That is because the backward momentum you get from the kick is transferred through friction with the ground you are standing on to the earth. The earth gets absorbs the "`backwards"' momentum, which is not measurable because of its huge mass compared to the ball's.}.\\

\noindent The amount of linear motion should somehow depend on how fast an object is moving. The faster it goes, the "`more motion"' there is. Also, the more "`object there is"'(i.e. the heavier it is) the more motion there is. As you are heavier than the ball, you will be moving backwards at a lesser speed than the ball is moving forward at. The speeds will be just right for the total amount of motion to be zero as you and the ball are moving in opposite directions.\\

\noindent According to this definition of momentum light does not have momentum at all because it does not have any mass\footnote{Some of the properties of light can be rather accurately described as if light was made of tiny mass-less particles called photons.}. From quantum mechanics we learn that mass-less particles do carry momentum and that the definition of momentum described above needs to be reviewed in such cases. Not only does light carry momentum but this momentum can be used to push other objects. Similarly to how one billiard ball can push another into motion, a photon can push another object into motion. Two phenomena that are instances of such transfer of momentum from light to other objects are
\begin{enumerate}
\item Compton effect: a photon colliding into an electron and pushing it into motion.
\item Solar sails: a kind of spacecraft made up of huge "`mirror sails"' that use light in a similar way to how a sailboat used wind. 
\end{enumerate}

\subsubsection{Angular momentum}

\noindent A similar phenomenon holds for angular momentum, the amount of rotational motion. When an ice-skater turns around his/her own axis, you often see them starting their rotation with their arms stretched out, away from their bodies. And during the turn they pull their arms inwards, towards their bodies. You see them turning much faster all of a sudden. This is an instance of conservation of angular momentum.\\

\noindent The amount of rotational motion there is will also depend on how much mass is rotating and on how fast it is rotating. But it also depends on how far away from the axis of rotation the rotating mass is. The further it is from its rotation axis, the more angular momentum it has. When the ice-skater rotates with their arms stretched out they have some amount of angular momentum. When pulling their arms inwards, the angular momentum is at risk of decreasing. As this is not allowed due to its conservation, it is compensated by increasing the speed of rotation, in such a way that the total angular momentum remains constant. 

\section{You are soooo cool...}
\label{Temperature_Intro}

\noindent In this section I want to consider a number of questions all related to temperature:
\begin{itemize}
\item What is temperature? What does it mean for an object to be hotter than another?
\item What happens when objects with different temperatures are placed close to each other?
\item How is an object's temperature related to our sensation of how warm or cold it is?
\item Is there an upper or a lower limit to the temperature of objects?
\item Do awesome and/or unexpected things happen at extreme temperatures?
\end{itemize}
The answer to the last question is yes. Things happen at extreme temperatures that is beyond most our imaginations. But to get there we will first consider the other questions more or less in the order presented.

\subsection{What is temperature?}

\noindent It is that quantity which is measured by a thermometer. In daily life it is expressed in degrees centigrade\footnote{Degrees Fahrenheit are usually used in the USA, at least outside of the scientific community.}. But what does a thermometer actually measure? Does an object with a higher temperature contain more (or less) of something?\\

\noindent Since Einstein's time we know that the material world around us is made of small constituents which are generically called particles. These particles are constantly moving and shaking around. Some faster than others. And as they are all moving around they have kinetic energy. Temperature turns out to be related to the average kinetic energy of the particles. Whenever you heat up some material, you are actually pumping energy into the material such that the particles have on average more kinetic energy (which for the moment we can interpret as "`the particles move on average at higher speeds"'\footnote{In section (\ref{} on relativity) we will see that this interpretation needs to be nuanced.}. An object with a higher temperature does indeed contain more of something than a colder object: energy.\\

\noindent From experience you probably know that two objects close to each always reach an equal temperature. Apparently the energy can "`flow"' out of one object (the hotter one) and into the other one (the colder one). The mechanism(s) by which this energy transfer occurs is(are) a topic to discuss in chapter \ref{}. However the energy "`moves"', it does move. Always from hot to cold. You may ask in that case how a refrigerator can work. We will also discuss this point in chapter \ref{} and see that also in this case the energy only flows from hot to cold at every step of the process, though it surely does not seem that way.\\

\noindent Another instance of a word meaning something else in science than it does in daily life is heat. We usually use the words heat and temperature more or less interchangeably. In physics they mean two very different things:
\begin{itemize}
\item \textbf{Temperature:} a quantity related to the amount of internal energy an object has. It is measured with a thermometer and expressed in decrees centigrade (another unit is usually used in the scientific context, but that is for later).
\item \textbf{Heat: }the amount of energy transferred from a hot to a cold object causing a change in temperature. The unit of heat is therefore that of energy and not degrees centigrade.
\end{itemize}

\noindent The fact that objects tend to reach a common temperature when close to each other is knows as the \textbf{zeroth law of thermodynamics}. Quite a fancy name for a seemingly simple fact\footnote{It is the zeroth one because there already was a first and second one, and this one was deemed too important to be the third one.}.

\subsection{The sensation of hotness}

\noindent While waiting for the train on a cold winters day find yourself a nice wooden bench. Touch it. Go on, touch it. Touch the wood. Then the metal. Did you feel it? Did you feel the difference?\\

\noindent That bench has probably been there for quite a while already and the zeroth law tells us that both the wood and the metal must have the same temperature. Otherwise they would exchange energy until they reach a common temperature. But to our touch they do not feel equally warm.\\

\noindent Both wood and metal do have the same temperature, which you can ascertain by attaching a thermometer to both of them. So why does it feel like the metal is much colder than the wood? It has to do with how good these materials are with transferring the heat. Your hand is warmer than the bench and there will be therefore be heat from your hand to the bench. But some materials (like most metals) are very good at absorbing that heat while others (like wood) are less so. We say that metal is a better \textbf{heat conductor} than wood. Our sensation of warm and cold is strongly related to the amount of energy that is "`sucked"' out of our hand every second. The better a material conducts heat, the higher the amount of energy it can absorb per second, the colder it seems to the touch.

\subsection{Limits of temperature}

\noindent Could we heat up an object indefinitely? In principle yes. An object would melt and later evaporate if we keep on heating it up. It would even do some funkier things after that like become a plasma. And we would not be able to keep it contained in anything, because the container itself at some point would not be able to withstand the temperature. Also, we would need to get that energy from somewhere and there is a finite amount of energy in the world. But there is as such, no reason known why it would not be possible to keep increasing the temperature of the object if an energy source and a mechanism to transfer the energy into the object are found.\\

\noindent Could we cool down an object indefinitely? No. There is a coldest temperature beyond which it is impossible to go. As you cool down an object you are depleting its internal energy, the molecules and atoms in the object are moving ever slower. At some point you will have depleted all its energy and all the constituents of the material will be standing still. It is simply impossible to make the particles move less than not moving. At that point the material has reached is lowest temperature possible.\\

\noindent There is at this point no reason to assume all materials will have the same lowest temperature. Iron, Carbon, oxygen and uranium have very different properties, so why would their lowest temperature be equal? Whatever the reason, the fact of the matter is all materials in the universe do have the same lowest temperature: $-273^{\circ} C$, absolute zero. Nothing in the universe can ever be colder nor can anything ever reach this temperature (quantum mechanics does not allow it), though we can get very close. The coldest parts of the universe outside of man-made labs is approximately $-270,45^{\circ} C$. This is a bit more than $2,5^{\circ} C$ warmer than absolute zero. This energy, small as it may be, is a remnant of the big bang, from a time that the universe was much warmer and denser.\\

\noindent To the best of my knowledge the coldest man-made temperature is a chilly $10^{-10}\; ^{\circ} C$ which is $0,000 \; 000 \; 000 \; 1^{\circ} C$. Quite close to absolute zero.\\

\noindent How do we get to such extremely low temperatures? Usually when we want to cool down object A, we bring object B which is colder with us and put it close to A. The former will heat up while the latter will cool down. For instance when we want a glass of coke to cool down, we put cold ice-blocks in it. But how can we cool down the coke below the temperature of the coldest ice-block we can find? The story is a bit more complicated than I describe here, but one of the techniques that was developed/used is called laser-cooling and it is basically exactly what the name suggests. You use a laser to cool things down. To understand we need to consider two facts that we've already discussed: 
\begin{enumerate}
\item temperature is a measure of kinetic energy.
\item light can transfer momentum to other objects.
\end{enumerate}
Laser light is shone on the material that is being cooled down. The light is made in such a way that its photons collide with the particles of the material in exactly the right way and the right moment to transfer momentum to the particles against their own momentum, thereby slowing them down. As they slow down their kinetic energy decreases and therefore so does their temperature.

\subsection{Cool phenomena}

\noindent Our intuition about how the world evolved in an environment in which the physical circumstances were "`normal"'. Our ancestors never directly experienced the very fast, the very small, the very warm, the very heavy, etc. Their world was confined to temperatures below $1000^{\circ} C$ and above $-70^{\circ} C$ and to objects never moving faster than a couple of hundreds of meters per second. Therefore our brains did not develop an intuition for what nature does outside of these ranges. Most of us will just assume that when objects move very fast or become extremely cold they will just be more extreme version of the things we already know. When we decrease the temperature to $-270^{\circ} C$, many will think it will just be a colder version of whatever happens at $-20^{\circ}C$. But when you think about it, there is no a priori reason to assume that is the case. As we know today, it often is not so at all. Which is why modern physical theories like quantum mechanics or relativity seem so difficult. They contradict what our intuition tells us should happen. They are demanding because we need to stretch our imagination and accept that our intuition is completely wrong for certain realms of nature.\\

\noindent I will briefly discuss three phenomena that occur at extremely low temperatures in some materials and that are unexpected and indeed counter-intuitive:
\begin{enumerate}
\item \textbf{Superconductivity: }Connect a battery to an electric wire made of copper and a current will flow through it. Connect the same battery to a wire made of gold and a weaker current will flow through it. Scientists say that gold has a higher electrical resistance than copper\footnote{It is actually the resistivity of copper that is lower than that of gold, not necessarily its resistance. There is a difference between both concepts, but for our current discussion it is very little importance.}. The resistance of materials like wood and plastic are much higher still and they are known as electrical insulators. Whether conductor or insulator, all materials have some electrical resistance. When a current flows through them, the material "`resists"' the electric flow and some of the electrical energy that went into it is "`lost"'. It leaves the wire in the form of heat. That is why electrical appliances heat up when used\footnote{Energy as you know is not lost, but only changes from "`useful"' electrical energy to "`useless"' heat. In most appliances the heat is not what we want, but for instance when ironing clothes, it is exactly this "`lost"' energy which is used.}.\\

\noindent Surprisingly, some materials (not necessarily conductors) have the special property that they completely lose their resistance when cooled down below some critical temperature. These are known as superconductors, though the lack of resistance is not their only peculiar property. When a superconductor is brought close to a magnet, it becomes a magnet itself and repels the original magnet. Therefore a magnet will levitate when put above a superconductor. The trains known as maglevs (for magnetic levitation) are based on this phenomenon called the Meissner effect.\\

\noindent Superconductivity was first discovered in solid mercury which was cooled down to approximately $-269^{\circ}C$ (which is $4^{\circ}C$ above absolute zero). A theory known as BCS-theory (after Bardeen, Cooper and Schrieffer who proposed it) was later developed to explain/describe the underlying, microscopic workings of it. The mechanism described in BCS only works for very low temperatures and according to it, no superconductive behavior should be present above around $-240^{\circ} C$.. But in the 1980's a "`high-temperature superconductor"' was discovered\footnote{High temperature is of course a relative concept, because we are still talking about $-180^{\circ} C$. This is way warmer than the original superconductors. Though still rather cold in my own humble opinion.}. The mechanism causing superconductivity at "`high"' temperatures is as yet unknown. It is possible that BCS is correct for low-temperature superconductivity and that a completely different mechanism causes its high-temperature counterpart. Alternatively, it could be wrong or at least incomplete even for the low-temperature ones. The question of the causes of high-temperature superconductivity remains open for now.

\item \textbf{Superfluidity: }Helium, known from balloons flying off and funny voices when inhaled has a number of surprising properties. Most elements in gaseous form liquify when cooled down below their boiling point and solidify when cooled down even further below their freezing point. Helium has the lowest boiling point of all known elements: $-269^{\circ} C$. But no matter how much it is cooled down further, it will never solidify under normal pressure conditions. This is a consequence of Heisenberg's quantum mechanical uncertainty principle\footnote{If you have not yet heard about this, do not worry. Keep on reading and you will in section (\ref{}).}. Though it does have some awesome properties when cooled down even further: it becomes \textbf{superfluid}. This is characterized by:
\begin{itemize}
\item \textbf{Zero viscosity: }stir a cup of coffee and leave it be. The speed with which the liquid rotates will decrease and after a short time stop altogether due to frictional forces between the different "`layers"' in the liquid. The amount of friction a fluid exerts on itself is called viscosity. Viscosity is a measure of how "`syrup-like"' a fluid is. Honey is much more viscous than water which is why it flows less fluently. Jam is even more viscous than honey.\\

\noindent  The more viscous a liquid is, the harder it is to make a thin layer of it. Take a tablespoon of oil and another of jam. Pour the oil on a tabletop and it will spread out. Do the same (on a clean part of the same tabletop) with jam and try to make it into a layer as thin as the oil. See what I mean? The less viscous a liquid is, the easier it is to make a thin layer of it.\\

\noindent When liquids are in contact with a solid (like the wall of the container in which they sit) they have the tendency to creep up or down the walls and "`wet"' it as shown in figure \ref{Capillar_fig} (By now I guess you are wondering what all of this has to do with cold liquid helium. Bear with me for just a short while longer. I promise, there is a point to this.). As the water tries to move up the container wall, it has to flow over itself, thereby experiencing viscous friction (usually referred to as viscous drag). It is the water's viscosity that prohibits it from going further up.\\

%\begin{figure}[h]
%\includegraphics[scale=0.2]{Capillarity}
%\centering
%\caption{Water creeps upwards in a thin glass tube and mercury creeps downwards}
%\label{Capillar_fig}
%\end{figure}
\noindent A superfluid has no viscosity whatsoever. There is therefore no viscous drag to terminate the "`crawl"' of the liquid up the walls of its container. It escapes from it all on its own. Awesome, I know!\\

\noindent As we are already on the topic of wetting, I want to mention capillarity. Have you ever wondered how plants and trees "`drink"'? They have roots of course, but how does the water get all the way to the top leaves? Have again a look at figure \ref{Capillar_fig}, more specifically at the thinner tube in the water. When a thin tube is put with one or both of its ends in a liquid, the liquid will be sucked into it. This phenomenon is called capillarity. Tree-trunks, branches and leaves all have thin tubes in them and it is capillarity that allows the water to rise from the roots in the ground all the way to the leaves, against the force of gravity.

\item \textbf{Quantized rotation: }when a glass of water is rotated the water itself will start rotating too as a consequence of the viscous drag of the inner walls of the glass on the water. As a superfluid is viscosity-free, it will not rotate together with its container. It will just keep standing still within a rotating container. 
\end{itemize}

\item \textbf{Bose-Einstein condensation: } Materials have three well-known phases: solid, liquid and gaseous. In the 1920's a young Indian physicist named Satyendra Nath Bose and Albert Einstein predicted the existence of a new phase (for some materials made of particles we now call bosons). We had to wait until 1995 for the first actual realization of this phase at the University of Colorado at Boulder and at MIT. The exact behavior of this phase is difficult to discuss without the necessary mathematics, which will come in a later section. I nevertheless want to emphasize here already how impressive it is to be able to predict a completely unknown natural phenomenon 70 years before it was ever realized.
\end{enumerate}